\documentstyle[12pt,fullpage,steve]{article}
\setcounter{page}{1}

\begin{document}

\title{EE 150 {\bf Bonus} Quiz \#8, {\large THIS QUIZ IS {\bf CLOSED} BOOK}}
\author{}
\date{}
\maketitle

ONLY tell what value is printed for each of the print statements...
\begin{verbatim}
#include <stdio.h>
main()
{
int a, b;
int c; float d, f;
int *p1, *p2;  
int e = 10; 
double t[3] = {2.0,2.1,2.2};
double *p3;
b = 2, a = b++;
printf("Question 1 a is: [ %i ]\n", a); 
d = 2.5, c = ++d;
printf("Question 2 a is : [ %i ]\n", a); 
printf("Question 3 expr is : [ %5.2f ]\n", ((int) (6.55 + 4.44)) * 3.1 );
printf("Question 4 expr is : [ %5.2f ]\n", ( (double) (6.55 + 4.44) * (int) 3.1 ) );
p1 = &e; p2 = p1, *p2 += 3;  *p1 = 11; a = *p1, *p2;  
printf("Question 5 a is: [ %i ]\n", a);
*(&a) = *t, p3 = t + 2, d = a + *p3;
printf("Question 6 d is : [ %5.2f ]\n", d);
printf("Question 7 a is : [ %i ]\n", a);
*(&(*(&d))) = *(t + 2), *(t + 2) = 3; f = t[2];
printf("Question 8: d is [ %5.2f ]\n", d);
printf("Question 9: f is [ %5.2f ]\n", f);
a = 5; b = 3; c = 1; c = (a == b) + (a > b) * (c < a) - (c < a > b);
printf("Question 10: c is [ %5.2f ]\n", c);
}
\end{verbatim}

My answers:
\begin{verbatim}
1:             2:           3:            4:             5:
6:             7:           8:            9:            10:
\end{verbatim}

\end{document}
